% =====================================================================
%  IEL04 - Circuitos Eléctricos I
%  Semana 7: Cálculo y medición de parámetros en circuitos RC en serie, en paralelo y mixtos: método, filtros, instrumentos y nomenclatura
%  Institución Universitaria de Barranquilla (IUB)
%  Generada por src/presentaciones.py (diseño IUB 5). Compilador: pdfLaTeX (dos pasadas)
% =====================================================================
\documentclass[aspectratio=169,11pt,xcolor={table}]{beamer}

% ---------------------------------------------------------------------
%  Codificación, idioma y tipografía
% ---------------------------------------------------------------------
\usepackage[utf8]{inputenc}
\usepackage[T1]{fontenc}
\usepackage[spanish,es-nodecimaldot,es-noshorthands]{babel}
\usepackage{avant}                       % Sans geométrica (emula Avenir Next)
\renewcommand{\familydefault}{\sfdefault}
\renewcommand{\ttdefault}{pcr}           % Monoespaciada para código
\usepackage{textcomp}

% ---------------------------------------------------------------------
%  Paquetes
% ---------------------------------------------------------------------
\usepackage{amsmath,amssymb,mathtools}
\usepackage{siunitx}
\sisetup{output-decimal-marker={.}, per-mode=symbol}
\usepackage{booktabs,tabularx,multirow,array}
\usepackage{adjustbox}
\usepackage{tikz}
\usetikzlibrary{arrows.meta,positioning,calc,shapes.geometric,shapes.misc,
  decorations.pathreplacing,fit,backgrounds,babel}
\usepackage[siunitx,RPvoltages]{circuitikz}
\usepackage{pgfplots}
\pgfplotsset{compat=1.18}
\usepackage[most]{tcolorbox}
\usepackage{listings}

% ---------------------------------------------------------------------
%  Paleta institucional IUB (Manual de Identidad Visual 2025)
% ---------------------------------------------------------------------
\definecolor{iubwhite}{HTML}{FFFFFF}
\definecolor{iubnavy}{HTML}{121023}
\definecolor{iubyellow}{HTML}{FFDF2D}
\definecolor{iubblue}{HTML}{00ADE7}
\definecolor{iuborange}{HTML}{E39037}
\definecolor{iubgreen}{HTML}{3B9C6D}

% ---------------------------------------------------------------------
%  Configuración de Beamer
% ---------------------------------------------------------------------
\setbeamertemplate{navigation symbols}{}
\setbeamersize{text margin left=0.6cm, text margin right=0.6cm}
\setbeamercolor{normal text}{fg=iubnavy, bg=iubwhite}
\setbeamercolor{background canvas}{bg=iubwhite}
\setbeamercolor{structure}{fg=iubnavy}
\setbeamercolor{alerted text}{fg=iuborange}
\setbeamertemplate{itemize item}{\textcolor{iubblue}{\small$\blacktriangleright$}}
\setbeamertemplate{itemize subitem}{\textcolor{iuborange}{\scriptsize$\bullet$}}
\setbeamertemplate{enumerate item}{\textcolor{iubblue}{\bfseries\insertenumlabel.}}
\setbeamertemplate{frametitle}{}         % el título vive en el headline

% Parches: el headline se pre-mide antes del primer frame
\makeatletter
\providecommand{\insertframetitle}{}
\providecommand{\insertframesubtitle}{}
\makeatother

% ---------- Encabezado ----------
\setbeamertemplate{headline}{%
  \begin{tikzpicture}
    \useasboundingbox (0,0) rectangle (\paperwidth,1.05cm);
    \fill[iubnavy] (0,0) rectangle (\paperwidth,1.05cm);
    \fill[iubyellow] (0,0) rectangle (\paperwidth,0.05cm);
    \node[anchor=west, text=iubyellow, font=\large\bfseries]
      at (0.6cm,0.55cm) {\strut\insertframetitle};
    \node[anchor=east, text=iubyellow, font=\footnotesize\bfseries]
      at ([xshift=-0.6cm]\paperwidth,0.55cm) {IUB};
  \end{tikzpicture}}

% ---------- Pie de página ----------
\setbeamertemplate{footline}{%
  \begin{tikzpicture}
    \useasboundingbox (0,0) rectangle (\paperwidth,0.55cm);
    \fill[iubnavy] (0,0) rectangle (\paperwidth,0.55cm);
    \node[anchor=west, text=iubyellow, font=\tiny]
      at (0.6cm,0.275cm) {IEL04 \textperiodcentered{} Circuitos Eléctricos I};
    \node[anchor=center, text=iubyellow, font=\tiny]
      at (0.66\paperwidth,0.275cm) {Ing. Sergio Martinez Castilla};
    \node[anchor=east, text=iubyellow, font=\tiny\bfseries]
      at ([xshift=-0.6cm]\paperwidth,0.275cm)
      {Semana 7 \quad \insertframenumber\,/\,\inserttotalframenumber};
  \end{tikzpicture}}

% ---------------------------------------------------------------------
%  Cajas tcolorbox (texto blanco sobre la paleta complementaria)
%  \begin{caja}[opciones]{color}{título} ... \end{caja}
%  \begin{cajaplana}[opciones]{color} ... \end{cajaplana}
%  \begin{cardB|cardO|cardG|cardN}[opciones]{título} ... \end{cardX}
% ---------------------------------------------------------------------
\tcbset{
  iubbase/.style={enhanced, arc=1.5mm, boxrule=0pt, coltext=white, coltitle=white,
    fonttitle=\bfseries\footnotesize, fontupper=\footnotesize,
    left=1.8mm, right=1.8mm, top=1mm, bottom=1.2mm,
    toptitle=0.6mm, bottomtitle=0.6mm, halign=left,
    before upper={\hyphenpenalty=5000\setbeamercolor{itemize item}{fg=white}%
                  \setbeamercolor{itemize/enumerate body}{fg=white}%
                  \setbeamercolor{itemize/enumerate subbody}{fg=white}%
                  \setbeamercolor{itemize subitem}{fg=white}%
                  \setbeamercolor{enumerate item}{fg=white}%
                  \setbeamertemplate{itemize item}{\textcolor{white}{\small$\blacktriangleright$}}}}
}
\newtcolorbox{caja}[3][]{iubbase, colback=#2, colframe=#2,
  colbacktitle=#2!78!iubnavy, title={#3}, #1}
\newtcolorbox{cajaplana}[2][]{iubbase, colback=#2, colframe=#2, #1}
\newtcolorbox{cardB}[2][]{iubbase, colback=iubblue,   colframe=iubblue,
  colbacktitle=iubblue!78!iubnavy,   title={#2}, #1}
\newtcolorbox{cardO}[2][]{iubbase, colback=iuborange, colframe=iuborange,
  colbacktitle=iuborange!78!iubnavy, title={#2}, #1}
\newtcolorbox{cardG}[2][]{iubbase, colback=iubgreen,  colframe=iubgreen,
  colbacktitle=iubgreen!78!iubnavy,  title={#2}, #1}
\newtcolorbox{cardN}[2][]{iubbase, colback=iubnavy,   colframe=iubnavy,
  colbacktitle=iubnavy, coltitle=iubyellow, title={#2}, #1}

% Celda de encabezado de tabla (texto amarillo sobre \rowcolor{iubnavy}) y código en línea
\newcommand{\hd}[1]{\textcolor{iubyellow}{\bfseries #1}}
\newcommand{\thc}[1]{\textcolor{iubyellow}{\textbf{#1}}}
\newcommand{\cod}[1]{\texttt{\textbf{#1}}}

% ---------------------------------------------------------------------
%  Código sobre fondo oscuro:
%  \begin{codebox}[listing options app={language=Python,firstnumber=1}]{título}
% ---------------------------------------------------------------------
\lstdefinestyle{iubcode}{
  basicstyle=\ttfamily\fontsize{7}{8.4}\selectfont\color{white},
  keywordstyle=\color{iubblue}\bfseries,
  commentstyle=\color{iubgreen!55!white}\itshape,
  stringstyle=\color{iuborange!80!white},
  numbers=left, numberstyle=\tiny\color{iubyellow}, numbersep=6pt,
  showstringspaces=false, breaklines=true, tabsize=4,
  columns=fullflexible, keepspaces=true, upquote=true}
\newtcblisting{codebox}[2][]{enhanced, listing only, colback=iubnavy,
  colframe=iubnavy, colbacktitle=iubnavy!85!white, coltitle=iubyellow,
  fonttitle=\bfseries\scriptsize, title={#2}, arc=1.5mm, boxrule=0pt,
  left=6mm, right=1mm, top=0.5mm, bottom=0.5mm, toptitle=0.5mm, bottomtitle=0.5mm,
  listing options={style=iubcode}, #1}

% ---------------------------------------------------------------------
%  Estilos TikZ
% ---------------------------------------------------------------------
\tikzset{
  bloque/.style={rectangle, rounded corners=2mm, fill=#1, text=white,
    font=\scriptsize\bfseries, align=center, minimum height=1.1cm,
    text width=1.8cm, inner sep=1.2mm},
  flecha/.style={-{Stealth[length=2.2mm]}, line width=1pt, draw=iubnavy},
  arr/.style={flecha},
  term/.style={rectangle, draw=#1, line width=1.2pt, fill=white,
    text=iubnavy, font=\tiny\bfseries, minimum width=1.05cm,
    minimum height=0.5cm, inner sep=1pt},
  nodo/.style={rectangle, rounded corners=1mm, fill=#1, text=white,
    font=\scriptsize\bfseries, minimum size=0.75cm, align=center},
  cable/.style={line width=1.4pt, draw=#1},
  box/.style={rectangle, rounded corners=1.5mm, draw=none, text=white,
    font=\scriptsize\bfseries, align=center, minimum height=0.8cm, inner sep=3pt},
  bB/.style={box, fill=iubblue}, bO/.style={box, fill=iuborange},
  bG/.style={box, fill=iubgreen}, bN/.style={box, fill=iubnavy},
  dec/.style={diamond, aspect=1.9, fill=iuborange, text=white,
    font=\scriptsize\bfseries, align=center, inner sep=1pt},
  tag/.style={fill=#1, text=white, font=\scriptsize\bfseries, rounded corners=1mm,
    minimum width=0.7cm, anchor=north west},
}

% ---------------------------------------------------------------------
%  Diapositiva separadora de bloque
%  #1 número  #2 título  #3 duración  #4 color
% ---------------------------------------------------------------------
\newcommand{\bloque}[4]{%
\section{#2}
\begin{frame}[plain]
\begin{tikzpicture}[remember picture, overlay]
  \fill[#4] (current page.north west) rectangle
        ([xshift=5.2cm]current page.south west);
  \node[anchor=north west, text=white, font=\bfseries\large]
    at ([xshift=0.8cm,yshift=-2.2cm]current page.north west) {BLOQUE};
  \node[anchor=north west, text=white, font=\bfseries\fontsize{72}{72}\selectfont]
    at ([xshift=0.65cm,yshift=-2.9cm]current page.north west) {#1};
  \node[anchor=north west, text=iubnavy, font=\bfseries\LARGE,
        text width=9.4cm, align=left]
    at ([xshift=6.2cm,yshift=-2.8cm]current page.north west)
    {\hyphenpenalty=10000\exhyphenpenalty=10000 #2};
  \node[anchor=north west, fill=iubnavy, text=iubyellow, rounded corners=1.5mm,
        font=\small\bfseries, inner sep=2mm]
    at ([xshift=6.2cm,yshift=-6.0cm]current page.north west) {Duración: #3};
  \fill[iubnavy] ([xshift=5.2cm]current page.south west) rectangle
        ([yshift=0.35cm]current page.south east);
  \fill[iubyellow] ([xshift=5.2cm,yshift=0.35cm]current page.south west)
        rectangle ([yshift=0.42cm]current page.south east);
  \node[anchor=north east, text=iubnavy, font=\small\bfseries]
    at ([xshift=-0.6cm,yshift=-0.5cm]current page.north east) {IUB · IEL04};
\end{tikzpicture}
\end{frame}}

% =====================================================================
\begin{document}
% =====================================================================

% ---------------------------------------------------------------------
%  PORTADA
% ---------------------------------------------------------------------
\begin{frame}[plain]
\begin{tikzpicture}[remember picture, overlay]
  % Panel institucional
  \fill[iubnavy] (current page.north west) rectangle
        ([xshift=5.6cm]current page.south west);
  \node[anchor=north west, text=iubyellow, font=\bfseries\fontsize{54}{54}\selectfont]
    at ([xshift=0.7cm,yshift=-1.0cm]current page.north west) {IUB};
  \node[anchor=north west, text=white, font=\footnotesize, text width=4.3cm, align=left]
    at ([xshift=0.75cm,yshift=-3.0cm]current page.north west)
    {Institución Universitaria de Barranquilla};
  \draw[iubyellow, line width=1.2pt]
    ([xshift=0.75cm,yshift=-4.25cm]current page.north west) --
    ([xshift=2.6cm,yshift=-4.25cm]current page.north west);
  \node[anchor=north west, text=iubyellow, font=\scriptsize\bfseries,
        text width=4.3cm, align=left]
    at ([xshift=0.75cm,yshift=-4.5cm]current page.north west)
    {Tecnología en Gestión de Sistemas Eléctricos};
  \node[anchor=south west, text=white, font=\scriptsize, text width=4.3cm, align=left]
    at ([xshift=0.75cm,yshift=0.9cm]current page.south west)
    {Ing. Sergio Martinez Castilla\\[1pt]\textcolor{iubyellow}{\tiny smartinezc@unibarranquilla.edu.co}};
  % Bloque de título
  \node[anchor=north west, fill=iubblue, text=white, rounded corners=1.5mm,
        font=\scriptsize\bfseries, inner sep=2mm]
    at ([xshift=6.4cm,yshift=-1.0cm]current page.north west)
    {IEL04 \textperiodcentered{} Semana 7 de 14 \textperiodcentered{} Corte evaluativo 2};
  \node[anchor=north west, text=iubnavy, font=\small, text width=9.0cm, align=left] (a)
    at ([xshift=6.4cm,yshift=-1.9cm]current page.north west)
    {Circuitos Eléctricos I};
  \node[anchor=north west, text=iubnavy, font=\bfseries\fontsize{13}{16}\selectfont,
        text width=9.0cm, align=left] (t)
    at ([yshift=-0.35cm]a.south west)
    {\hyphenpenalty=10000 Cálculo y medición de parámetros en circuitos RC en serie, en paralelo y mixtos: método, filtros, instrumentos y nomenclatura};
  % Franjas de la paleta complementaria
  \fill[iubblue]   ([xshift=6.4cm,yshift=1.1cm]current page.south west)
    rectangle ++(1.6cm,0.18cm);
  \fill[iuborange] ([xshift=8.1cm,yshift=1.1cm]current page.south west)
    rectangle ++(1.6cm,0.18cm);
  \fill[iubgreen]  ([xshift=9.8cm,yshift=1.1cm]current page.south west)
    rectangle ++(1.6cm,0.18cm);
  \node[anchor=north west, text=iubnavy, font=\scriptsize]
    at ([xshift=6.4cm,yshift=0.95cm]current page.south west)
    {Sesión teórico-práctica \textperiodcentered{} 240 min};
\end{tikzpicture}
\end{frame}

% ---------------------------------------------------------------------
\begin{frame}{Punto de partida de la sesión}
\begin{columns}[T,onlytextwidth]
  \column{0.34\textwidth}
\begin{caja}[height=5.9cm]{iubblue}{Continuidad}
\scriptsize \begin{itemize}\setlength{\itemsep}{2pt}
  \item Semana 6: Circuito RC en paralelo
  \item \textbf{Semana 7: Cálculo y medición de parámetros en circuitos RC serie y paralelo: unidades, simbología y nomenclatura}
  \item Semana 8: Circuitos RL en serie y en paralelo
\end{itemize}
\end{caja}
  \column{0.34\textwidth}
\begin{caja}[height=5.9cm]{iuborange}{Prerrequisitos}
\scriptsize \begin{itemize}\setlength{\itemsep}{2pt}
  \item Calcular impedancia, ángulo de fase y tensiones de un RC en serie
  \item Calcular admitancia y corrientes de un RC en paralelo aplicando la LCK
  \item Convertir fasores entre forma rectangular y polar, y operar con ellos
  \item Identificar y verificar resistores y condensadores con ohmímetro
\end{itemize}
\end{caja}
  \column{0.29\textwidth}
\begin{caja}[height=5.9cm]{iubgreen}{Competencia}
\scriptsize \textbf{UC2.} Coordinar actividades asociadas a sistemas eléctricos utilizando fuentes convencionales y no convencionales de energía.
\end{caja}
\end{columns}
\end{frame}

% ---------------------------------------------------------------------
\begin{frame}{Resultados de aprendizaje}
\begin{columns}[c,onlytextwidth]
  \column{0.45\textwidth}
\begin{caja}{iubnavy}{IEL04-RA2}
\footnotesize Calcular un circuito resistivo, capacitivo y RC, sea en serie o paralelo.
\end{caja}
\vspace{1mm}
\begin{cajaplana}{iubblue}
\scriptsize \textbf{CE1.} Identificar los tipos de resistores y capacitores.\\[2pt]
\textbf{CE2.} Realizar mediciones en un circuito resistivo y capacitivo.\\[2pt]
\textbf{CE3.} Realizar mediciones en un circuito RC.
\end{cajaplana}
  \column{0.52\textwidth}
\centering
\begin{adjustbox}{max width=\linewidth, max totalheight=6.1cm}
\begin{tikzpicture}
  \node[box, fill=iubblue, text width=4.9cm, align=left, font=\tiny, anchor=west] (r0) at (1.25,0.00) {Identificar los componentes y la topología de un circuito RC y elegir el método de cálculo adecuado para cada parte.};
  \node[tag=iubnavy, font=\tiny\bfseries, minimum width=1.15cm, anchor=east] at (1.1,0.00) {Analizar};
  \node[box, fill=iuborange, text width=4.9cm, align=left, font=\tiny, anchor=west] (r1) at (1.25,-1.45) {Calcular impedancias, corrientes y tensiones de circuitos RC en serie, en paralelo y mixtos con fasores.};
  \node[tag=iubnavy, font=\tiny\bfseries, minimum width=1.15cm, anchor=east] at (1.1,-1.45) {Aplicar};
  \node[box, fill=iubgreen, text width=4.9cm, align=left, font=\tiny, anchor=west] (r2) at (1.25,-2.90) {Calcular la frecuencia de corte de un filtro RC y predecir la tensión de salida a una frecuencia dada.};
  \node[tag=iubnavy, font=\tiny\bfseries, minimum width=1.15cm, anchor=east] at (1.1,-2.90) {Aplicar};
  \node[box, fill=iubblue, text width=4.9cm, align=left, font=\tiny, anchor=west] (r3) at (1.25,-4.35) {Medir tensiones y corrientes en un circuito RC mixto, evaluar el efecto de los instrumentos y justificar las diferencias con el cálculo.};
  \node[tag=iubnavy, font=\tiny\bfseries, minimum width=1.15cm, anchor=east] at (1.1,-4.35) {Evaluar};
  \draw[flecha, draw=iubnavy!40] (r0.south) -- (r1.north);
  \draw[flecha, draw=iubnavy!40] (r1.south) -- (r2.north);
  \draw[flecha, draw=iubnavy!40] (r2.south) -- (r3.north);
\end{tikzpicture}
\end{adjustbox}
\end{columns}
\end{frame}

% ---------------------------------------------------------------------
\begin{frame}{Ruta de la sesión \textperiodcentered{} 240 minutos}
\centering
\begin{tikzpicture}[x=1cm,y=1cm]
  \node[circle, fill=iubblue, text=white, font=\scriptsize\bfseries, minimum size=0.6cm, inner sep=0] at (0,0.00) {1};
  \node[anchor=west, font=\scriptsize, text width=7.6cm] at (0.45,0.00) {Integrar lo visto en serie y en paralelo en un procedimiento de cálculo y medición};
  \fill[iubblue] (8.3,-0.20) rectangle ++(1.84,0.4);
  \node[anchor=west, font=\scriptsize\bfseries] at (10.24,0.00) {20 min};
  \node[circle, fill=iuborange, text=white, font=\scriptsize\bfseries, minimum size=0.6cm, inner sep=0] at (0,-0.76) {2};
  \node[anchor=west, font=\scriptsize, text width=7.6cm] at (0.45,-0.76) {Pasos del análisis fasorial: reactancias, reducción, ley de Ohm y verificación con Kirchhoff};
  \fill[iuborange] (8.3,-0.96) rectangle ++(2.76,0.4);
  \node[anchor=west, font=\scriptsize\bfseries] at (11.16,-0.76) {30 min};
  \node[circle, fill=iubgreen, text=white, font=\scriptsize\bfseries, minimum size=0.6cm, inner sep=0] at (0,-1.51) {3};
  \node[anchor=west, font=\scriptsize, text width=7.6cm] at (0.45,-1.51) {Reducción de un resistor en serie con un paralelo RC y reparto de tensiones y corrientes};
  \fill[iubgreen] (8.3,-1.71) rectangle ++(3.22,0.4);
  \node[anchor=west, font=\scriptsize\bfseries] at (11.62,-1.51) {35 min};
  \node[circle, fill=iubblue, text=white, font=\scriptsize\bfseries, minimum size=0.6cm, inner sep=0] at (0,-2.27) {4};
  \node[anchor=west, font=\scriptsize, text width=7.6cm] at (0.45,-2.27) {Filtros pasabajas y pasaaltas, fc = 1/(2\ensuremath{\pi}RC) y el 70.7 \%};
  \fill[iubblue] (8.3,-2.47) rectangle ++(3.22,0.4);
  \node[anchor=west, font=\scriptsize\bfseries] at (11.62,-2.27) {35 min};
  \node[circle, fill=iuborange, text=white, font=\scriptsize\bfseries, minimum size=0.6cm, inner sep=0] at (0,-3.03) {5};
  \node[anchor=west, font=\scriptsize, text width=7.6cm] at (0.45,-3.03) {Resistencia interna del generador, carga de la punta del osciloscopio, ancho de banda del multímetro};
  \fill[iuborange] (8.3,-3.23) rectangle ++(2.76,0.4);
  \node[anchor=west, font=\scriptsize\bfseries] at (11.16,-3.03) {30 min};
  \node[circle, fill=iubgreen, text=white, font=\scriptsize\bfseries, minimum size=0.6cm, inner sep=0] at (0,-3.79) {6};
  \node[anchor=west, font=\scriptsize, text width=7.6cm] at (0.45,-3.79) {Cuadro de magnitudes, unidades, prefijos, símbolos y notación fasorial};
  \fill[iubgreen] (8.3,-3.99) rectangle ++(2.30,0.4);
  \node[anchor=west, font=\scriptsize\bfseries] at (10.70,-3.79) {25 min};
  \node[circle, fill=iubblue, text=white, font=\scriptsize\bfseries, minimum size=0.6cm, inner sep=0] at (0,-4.54) {7};
  \node[anchor=west, font=\scriptsize, text width=7.6cm] at (0.45,-4.54) {Resistor en serie con un paralelo RC a 1 kHz: predecir, medir y justificar las diferencias};
  \fill[iubblue] (8.3,-4.74) rectangle ++(4.60,0.4);
  \node[anchor=west, font=\scriptsize\bfseries] at (13.00,-4.54) {50 min};
  \node[circle, fill=iuborange, text=white, font=\scriptsize\bfseries, minimum size=0.6cm, inner sep=0] at (0,-5.30) {8};
  \node[anchor=west, font=\scriptsize, text width=7.6cm] at (0.45,-5.30) {Síntesis, cierre actitudinal y bibliografía};
  \fill[iuborange] (8.3,-5.50) rectangle ++(1.38,0.4);
  \node[anchor=west, font=\scriptsize\bfseries] at (9.78,-5.30) {15 min};
  \draw[iubnavy, line width=0.6pt] (8.3,0.45) -- (8.3,-5.70);
\end{tikzpicture}
\end{frame}

% =====================================================================
\bloque{01}{Integrar lo visto en serie y en paralelo en un procedimiento de cálculo y medición}{20 min}{iubblue}
% =====================================================================

% ---------------------------------------------------------------------
\begin{frame}{Un solo método para todos los circuitos RC}
\begin{columns}[c,onlytextwidth]
  \column{0.57\textwidth}
\small
\begin{itemize}\setlength{\itemsep}{4pt}
  \item Pero los circuitos reales rara vez son puramente de un tipo.
\end{itemize}
  \column{0.4\textwidth}
\begin{cardO}{Nota pedagógica}
\footnotesize Esta sesión es también una preparación para el RA3: en la semana 8, el mismo diagrama sirve para los circuitos RL, cambiando solo el cálculo de la reactancia y el signo del ángulo.
\end{cardO}
\end{columns}
\end{frame}

% ---------------------------------------------------------------------
\begin{frame}{{\normalsize Procedimiento general de cálculo y medición de un circuito}}
\centering
\begin{adjustbox}{max width=\linewidth, max totalheight=6.1cm}
\begin{tikzpicture}
  \node[bG, align=center, font=\scriptsize, text width=3.51cm] (n0) at (1.88,4.53) {Identificar y verificar componentes};
  \node[bB, align=center, font=\scriptsize, text width=2.80cm] (n1) at (6.77,4.53) {XC a la frecuencia de trabajo};
  \node[bB, align=center, font=\scriptsize, text width=3.03cm] (n2) at (11.43,4.53) {Reducir: Z en serie, Y en paralelo};
  \node[bB, align=center, font=\scriptsize, text width=3.03cm] (n3) at (11.43,1.07) {Ley de Ohm fasorial: I y V};
  \node[bO, align=center, font=\scriptsize, text width=2.58cm] (n4) at (6.77,1.07) {Verificar LVK y LCK};
  \node[bG, align=center, font=\scriptsize, text width=2.84cm] (n5) at (1.88,1.07) {Medir y comparar};
  \draw[flecha] (n0) -- (n1);
  \draw[flecha] (n1) -- (n2);
  \draw[flecha] (n2) -- (n3);
  \draw[flecha] (n3) -- (n4);
  \draw[flecha] (n4) -- (n5);
\end{tikzpicture}
\end{adjustbox}

\vspace{1mm}{\scriptsize Procedimiento general de cálculo y medición de un circuito RC\par}
\end{frame}

% =====================================================================
\bloque{02}{Pasos del análisis fasorial: reactancias, reducción, ley de Ohm y verificación con Kirchhoff}{30 min}{iuborange}
% =====================================================================

% ---------------------------------------------------------------------
\begin{frame}{{\normalsize Método general de cálculo con impedancias y admitancias}}
\begin{columns}[c,onlytextwidth]
  \column{0.57\textwidth}
\small
\begin{itemize}\setlength{\itemsep}{4pt}
  \item Con estas herramientas, el método general tiene seis pasos.
  \item Cambiar la referencia gira todo el diagrama fasorial, pero no altera ni las magnitudes ni las diferencias de fase entre las magnitudes.
  \item En los circuitos mixtos se combinan ambos, como muestra la siguiente unidad.
\end{itemize}
  \column{0.4\textwidth}
\begin{caja}{iubblue}{Ecuación clave}
\footnotesize \centering\small \begin{adjustbox}{max width=\linewidth}$\displaystyle \begin{gathered}\mathbf{Z} = R + jX = Z\angle\theta \\ Z = \sqrt{R^{2} + X^{2}},\ \ \theta = \tan^{-1}\!\left(\frac{X}{R}\right) \\ R = Z\cos\theta,\ \ X = Z\sin\theta\end{gathered}$\end{adjustbox}
\end{caja}
\end{columns}
\end{frame}

% ---------------------------------------------------------------------
\begin{frame}{{\normalsize Método general aplicado al RC en serie de la semana 4}}
\centering\scriptsize
\renewcommand{\arraystretch}{1.35}
\rowcolors{2}{iubnavy!6}{white}
\begin{adjustbox}{max width=\linewidth, max totalheight=5.6cm}
\begin{tabularx}{\textwidth}{>{\raggedright\arraybackslash}X>{\raggedright\arraybackslash}X>{\raggedright\arraybackslash}X>{\raggedright\arraybackslash}X}
  \rowcolor{iubnavy}
  \hd{Paso} & \hd{Operación} & \hd{Forma} & \hd{Resultado}\\
  1 & Identificar y medir componentes & — & R = 1.5 k\ensuremath{\Omega}, C = 0.1 µF\\
  2 & XC = 1/(2\ensuremath{\pi}fC) & magnitud & 1592 \ensuremath{\Omega}\\
  3 & Z = R \ensuremath{-} jXC & rectangular \ensuremath{\rightarrow} polar & 2187 \ensuremath{\angle} \ensuremath{-}46.7° \ensuremath{\Omega}\\
  4 & I = V/Z & polar & 2.29 \ensuremath{\angle} 46.7° mA\\
  5 & VR = IR, VC = IXC & polar & 3.43 \ensuremath{\angle} 46.7° V; 3.64 \ensuremath{\angle} \ensuremath{-}43.3° V\\
  6 & VR + VC = VS & rectangular & 5.00 \ensuremath{\angle} 0° V\\
\end{tabularx}
\end{adjustbox}
\vspace{2mm}
{\scriptsize Método general aplicado al RC en serie de la semana 4 (1.5 k\ensuremath{\Omega}, 0.1 µF, 5 V eficaces a 1 kHz)\par}
\end{frame}

% =====================================================================
\bloque{03}{Reducción de un resistor en serie con un paralelo RC y reparto de tensiones y corrientes}{35 min}{iubgreen}
% =====================================================================

% ---------------------------------------------------------------------
\begin{frame}{Circuitos RC mixtos (serie-paralelo)}
\begin{columns}[c,onlytextwidth]
  \column{0.57\textwidth}
\small
\begin{itemize}\setlength{\itemsep}{4pt}
  \item Un circuito RC mixto combina grupos en serie y grupos en paralelo.
  \item Su admitancia es $\mathbf{Y}_p = G_2 + jB_C$, y su impedancia, $\mathbf{Z}_p = 1/\mathbf{Y}_p$, se pasa a forma rectangular para poder sumarla con $R_1$.
  \item Con 5 V eficaces a 0°, la corriente total es $\mathbf{I} = \mathbf{V}_S/\mathbf{Z}_T = 2.57 \angle 22.6^\circ\ \text{mA}$.
\end{itemize}
  \column{0.4\textwidth}
\begin{caja}{iubblue}{Ecuación clave}
\footnotesize \centering\small \begin{adjustbox}{max width=\linewidth}$\displaystyle \mathbf{Z}_T = R_1 + \mathbf{Z}_p = R_1 + \frac{1}{G_2 + jB_C}$\end{adjustbox}
\end{caja}
\end{columns}
\end{frame}

% ---------------------------------------------------------------------
\begin{frame}{Resultados del circuito mixto R1 + (R2 \ensuremath{\parallel} C) (1/2)}
\centering\scriptsize
\renewcommand{\arraystretch}{1.35}
\rowcolors{2}{iubnavy!6}{white}
\begin{adjustbox}{max width=\linewidth, max totalheight=5.6cm}
\begin{tabularx}{\textwidth}{>{\raggedright\arraybackslash}X>{\raggedright\arraybackslash}X>{\raggedright\arraybackslash}X>{\raggedright\arraybackslash}X}
  \rowcolor{iubnavy}
  \hd{Magnitud} & \hd{Valor} & \hd{Ángulo (ref. VS)} & \hd{Comprobación}\\
  ZT & 1945 \ensuremath{\Omega} & \ensuremath{-}22.6° & 1795 \ensuremath{-} j749 \ensuremath{\Omega}\\
  I & 2.57 mA & +22.6° & VS/ZT\\
  VR1 & 2.57 V & +22.6° & en fase con I\\
  Vp & 2.81 V & \ensuremath{-}20.7° & común a R2 y C\\
\end{tabularx}
\end{adjustbox}
\vspace{2mm}
{\scriptsize Resultados del circuito mixto R1 + (R2 \ensuremath{\parallel} C) con valores nominales (5 V eficaces a 1 kHz)\par}
\end{frame}

% ---------------------------------------------------------------------
\begin{frame}{Resultados del circuito mixto R1 + (R2 \ensuremath{\parallel} C) (2/2)}
\centering\scriptsize
\renewcommand{\arraystretch}{1.35}
\rowcolors{2}{iubnavy!6}{white}
\begin{adjustbox}{max width=\linewidth, max totalheight=5.6cm}
\begin{tabularx}{\textwidth}{>{\raggedright\arraybackslash}X>{\raggedright\arraybackslash}X>{\raggedright\arraybackslash}X>{\raggedright\arraybackslash}X}
  \rowcolor{iubnavy}
  \hd{Magnitud} & \hd{Valor} & \hd{Ángulo (ref. VS)} & \hd{Comprobación}\\
  IR2 & 1.87 mA & \ensuremath{-}20.7° & en fase con Vp\\
  IC & 1.76 mA & +69.3° & 90° delante de Vp\\
  \ensuremath{\surd}(IR2² + IC²) & 2.57 mA & — & LCK: igual a I\\
  VR1 + Vp (fasorial) & 5.00 V & 0° & LVK: igual a VS\\
\end{tabularx}
\end{adjustbox}
\vspace{2mm}
{\scriptsize Resultados del circuito mixto R1 + (R2 \ensuremath{\parallel} C) con valores nominales (5 V eficaces a 1 kHz)\par}
\end{frame}

% =====================================================================
\bloque{04}{Filtros pasabajas y pasaaltas, fc = 1/(2\ensuremath{\pi}RC) y el 70.7 \%}{35 min}{iubblue}
% =====================================================================

% ---------------------------------------------------------------------
\begin{frame}{El circuito RC como filtro}
\begin{columns}[c,onlytextwidth]
  \column{0.57\textwidth}
\small
\begin{itemize}\setlength{\itemsep}{4pt}
  \item La aplicación más extendida del circuito RC en serie es el \textbf{filtro}: un circuito que deja pasar unas frecuencias y atenúa otras.
  \item La frontera entre las frecuencias que pasan y las que se atenúan se define por convención.
  \item En el laboratorio, $f_c$ se mide barriendo la frecuencia del generador hasta que la salida cae al 70.7 \% de su valor en baja frecuencia.
\end{itemize}
  \column{0.4\textwidth}
\begin{caja}{iubblue}{Ecuación clave}
\footnotesize \centering\small \begin{adjustbox}{max width=\linewidth}$\displaystyle \begin{gathered}f_c = \frac{1}{2\pi R C} \\ \frac{V_{sal}}{V_{ent}}\bigg|_{pasabajas} = \frac{1}{\sqrt{1 + (f/f_c)^{2}}} \\ \frac{V_{sal}}{V_{ent}}\bigg|_{pasaaltas} = \frac{f/f_c}{\sqrt{1 + (f/f_c)^{2}}}\end{gathered}$\end{adjustbox}
\end{caja}
\end{columns}
\end{frame}

% ---------------------------------------------------------------------
\begin{frame}{Respuesta en frecuencia de un filtro RC pasabajas}
\begin{columns}[c,onlytextwidth]
  \column{0.62\textwidth}
\centering
\begin{tikzpicture}
  \begin{semilogxaxis}[width=8.4cm, height=5.7cm, xmin=10, xmax=1e+05, ymin=0, ymax=1.0594,
      xlabel={Frecuencia f (Hz)}, ylabel={Vsal / Vent},
      label style={font=\scriptsize, text=iubnavy}, tick label style={font=\tiny, text=iubnavy},
      axis line style={iubnavy}, grid=major, grid style={iubnavy!12},
      legend style={font=\tiny, fill=white}, legend pos=north east]
    \addplot[line width=1.4pt, iubblue] coordinates {(10,0.99996) (11.311,0.99994) (12.795,0.99993) (14.472,0.99991) (16.37,0.99988) (18.516,0.99985) (20.945,0.99981) (23.691,0.99975) (26.797,0.99968) (30.311,0.99959) (34.286,0.99948) (38.782,0.99933) (43.867,0.99915) (49.619,0.99891) (56.126,0.9986) (63.486,0.99821) (71.81,0.99772) (81.227,0.99708) (91.878,0.99627) (103.93,0.99524) (117.55,0.99392) (132.97,0.99224) (150.4,0.9901) (170.13,0.98739) (192.43,0.98395) (217.67,0.9796) (246.21,0.97412) (278.49,0.96724) (315.01,0.95864) (356.32,0.94797) (403.04,0.93482) (455.89,0.91878) (515.67,0.8994) (583.29,0.87631) (659.78,0.8492) (746.29,0.81793) (844.15,0.78254) (954.85,0.74331) (1080.1,0.70079) (1221.7,0.65571) (1381.9,0.609) (1563.1,0.56163) (1768,0.51456) (1999.9,0.46866) (2262.1,0.42464) (2558.7,0.38303) (2894.3,0.34419) (3273.8,0.3083) (3703.1,0.27544) (4188.6,0.24555) (4737.9,0.21853) (5359.2,0.19421) (6061.9,0.17241) (6856.8,0.15292) (7755.9,0.13554) (8772.9,0.12007) (9923.3,0.10631) (11225,0.094106) (12696,0.083277) (14361,0.073679) (16244,0.065176) (18374,0.057647) (20784,0.050983) (23509,0.045085) (26592,0.039868) (30079,0.035252) (34023,0.03117) (38484,0.027559) (43531,0.024366) (49239,0.021543) (55695,0.019047) (62999,0.016839) (71260,0.014888) (80604,0.013162) (91173,0.011636) (1e+05,0.010609)};
    \addlegendentry{Pasabajas (salida en C)}
    \addplot[line width=1.4pt, iuborange] coordinates {(10,0.0094247) (11.311,0.01066) (12.795,0.012058) (14.472,0.013639) (16.37,0.015427) (18.516,0.017449) (20.945,0.019737) (23.691,0.022323) (26.797,0.025249) (30.311,0.028557) (34.286,0.032298) (38.782,0.036528) (43.867,0.04131) (49.619,0.046716) (56.126,0.052825) (63.486,0.059729) (71.81,0.067527) (81.227,0.076333) (91.878,0.086273) (103.93,0.097484) (117.55,0.11012) (132.97,0.12435) (150.4,0.14035) (170.13,0.15832) (192.43,0.17846) (217.67,0.20097) (246.21,0.22605) (278.49,0.25388) (315.01,0.28462) (356.32,0.31836) (403.04,0.35511) (455.89,0.39478) (515.67,0.43713) (583.29,0.48176) (659.78,0.52807) (746.29,0.57532) (844.15,0.6226) (954.85,0.66894) (1080.1,0.71337) (1221.7,0.75501) (1381.9,0.79317) (1563.1,0.82739) (1768,0.85746) (1999.9,0.88338) (2262.1,0.90536) (2558.7,0.92373) (2894.3,0.9389) (3273.8,0.95129) (3703.1,0.96132) (4188.6,0.96938) (4737.9,0.97583) (5359.2,0.98096) (6061.9,0.98503) (6856.8,0.98824) (7755.9,0.99077) (8772.9,0.99277) (9923.3,0.99433) (11225,0.99556) (12696,0.99653) (14361,0.99728) (16244,0.99787) (18374,0.99834) (20784,0.9987) (23509,0.99898) (26592,0.9992) (30079,0.99938) (34023,0.99951) (38484,0.99962) (43531,0.9997) (49239,0.99977) (55695,0.99982) (62999,0.99986) (71260,0.99989) (80604,0.99991) (91173,0.99993) (1e+05,0.99994)};
    \addlegendentry{Pasaaltas (salida en R)}
    \addplot[line width=1pt, dashed, iubnavy!70] coordinates {(1061,0) (1061,1.0594)};
    \addlegendentry{fc = 1061 Hz}
    \addplot[line width=1pt, dashed, iubnavy!70] coordinates {(10,0.707) (1e+05,0.707)};
    \addlegendentry{70.7 \%}
  \end{semilogxaxis}
\end{tikzpicture}
  \column{0.35\textwidth}
\begin{caja}{iubblue}{Lectura de la gráfica}
\footnotesize La gráfica muestra que las dos curvas se cruzan en la frecuencia de corte, a 0.707, y son simétricas respecto a ella en escala logarítmica.
\end{caja}
\end{columns}
\end{frame}

% =====================================================================
\bloque{05}{Resistencia interna del generador, carga de la punta del osciloscopio, ancho de banda del multímetro}{30 min}{iuborange}
% =====================================================================

% ---------------------------------------------------------------------
\begin{frame}{{\normalsize Instrumentos de medición y sus efectos en circuitos RC}}
\begin{columns}[c,onlytextwidth]
  \column{0.57\textwidth}
\small
\begin{itemize}\setlength{\itemsep}{4pt}
  \item Todo instrumento que se conecta a un circuito se convierte, durante la medición, en parte del circuito.
  \item La \textbf{punta de osciloscopio} es el caso típico.
\end{itemize}
  \column{0.4\textwidth}
\begin{caja}{iubblue}{Ecuación clave}
\footnotesize \centering\small \begin{adjustbox}{max width=\linewidth}$\displaystyle X_{Cp} = \frac{1}{2\pi f C_p}$\end{adjustbox}
\end{caja}
\end{columns}
\end{frame}

% ---------------------------------------------------------------------
\begin{frame}{{\normalsize Efectos de los instrumentos en las mediciones de circuitos}}
\centering\scriptsize
\renewcommand{\arraystretch}{1.35}
\rowcolors{2}{iubnavy!6}{white}
\begin{adjustbox}{max width=\linewidth, max totalheight=5.6cm}
\begin{tabularx}{\textwidth}{>{\raggedright\arraybackslash}X>{\raggedright\arraybackslash}X>{\raggedright\arraybackslash}X>{\raggedright\arraybackslash}X}
  \rowcolor{iubnavy}
  \hd{Instrumento} & \hd{Parámetro interno} & \hd{Efecto sobre la medición} & \hd{Cómo reducirlo}\\
  Generador de funciones & resistencia de salida (\ensuremath{\approx}50 \ensuremath{\Omega}) & la tensión aplicada baja al conectar la carga & medir VS en los bornes del circuito\\
  Osciloscopio con punta ×1 & \ensuremath{\approx}1 M\ensuremath{\Omega} \ensuremath{\parallel} capacitancia alta & carga capacitiva a frecuencias altas & usar punta ×10 compensada\\
  Osciloscopio con punta ×10 & \ensuremath{\approx}10 M\ensuremath{\Omega} \ensuremath{\parallel} pocos pF & efecto mínimo & compensarla con la onda cuadrada de referencia\\
  Multímetro como voltímetro ca & ancho de banda limitado & lectura baja a frecuencias altas & verificar la especificación; usar el osciloscopio\\
  Multímetro como amperímetro & resistencia en serie & reduce ligeramente la corriente de la rama & preferir un resistor sensor pequeño\\
  Puntas con tierra común & tierras unidas al chasis & cortocircuita elementos flotantes & medir con resta de canales\\
\end{tabularx}
\end{adjustbox}
\vspace{2mm}
{\scriptsize Efectos de los instrumentos en las mediciones de circuitos RC (puntas ×1 y ×10: $[1]$; generador, multímetro y tierras: [fuente no cargada])\par}
\end{frame}

% =====================================================================
\bloque{06}{Cuadro de magnitudes, unidades, prefijos, símbolos y notación fasorial}{25 min}{iubgreen}
% =====================================================================

% ---------------------------------------------------------------------
\begin{frame}{Parámetros, unidades, simbología y nomenclatura}
\begin{columns}[c,onlytextwidth]
  \column{0.57\textwidth}
\small
\begin{itemize}\setlength{\itemsep}{4pt}
  \item El contenido procedimental del RA2 pide diferenciar los parámetros del circuito RC identificando unidades de medida, múltiplos y submúltiplos, simbología y nomenclatura.
  \item Tres observaciones evitan los errores más comunes al leer la tabla.
\end{itemize}
  \column{0.4\textwidth}
\begin{cardO}{Error frecuente}
\footnotesize Mezclar valores eficaces con valores pico a pico en un mismo cálculo.
\end{cardO}
\end{columns}
\end{frame}

% ---------------------------------------------------------------------
\begin{frame}{Parámetros de los circuitos RC (1/2)}
\centering\scriptsize
\renewcommand{\arraystretch}{1.35}
\rowcolors{2}{iubnavy!6}{white}
\begin{adjustbox}{max width=\linewidth, max totalheight=5.6cm}
\begin{tabularx}{\textwidth}{>{\raggedright\arraybackslash}X>{\raggedright\arraybackslash}X>{\raggedright\arraybackslash}X>{\raggedright\arraybackslash}X>{\raggedright\arraybackslash}X}
  \rowcolor{iubnavy}
  \hd{Parámetro} & \hd{Símbolo} & \hd{Unidad} & \hd{Submúltiplos habituales} & \hd{Fórmula o definición}\\
  Resistencia & R & ohmio (\ensuremath{\Omega}) & k\ensuremath{\Omega}, M\ensuremath{\Omega} & V/I en el resistor\\
  Capacitancia & C & faradio (F) & pF, nF, µF & Q/V\\
  Reactancia capacitiva & XC & ohmio (\ensuremath{\Omega}) & k\ensuremath{\Omega} & 1/(2\ensuremath{\pi}fC)\\
  Impedancia & Z, Z\ensuremath{\angle}\ensuremath{\theta} & ohmio (\ensuremath{\Omega}) & k\ensuremath{\Omega} & R \ensuremath{-} jXC (serie)\\
  Conductancia & G & siemens (S) & mS, µS & 1/R\\
  Susceptancia capacitiva & BC & siemens (S) & mS, µS & 2\ensuremath{\pi}fC\\
  Admitancia & Y, Y\ensuremath{\angle}\ensuremath{\theta} & siemens (S) & mS, µS & G + jBC (paralelo)\\
\end{tabularx}
\end{adjustbox}
\vspace{2mm}
{\scriptsize Parámetros de los circuitos RC: símbolo, unidad, submúltiplos habituales y fórmula\par}
\end{frame}

% ---------------------------------------------------------------------
\begin{frame}{Parámetros de los circuitos RC (2/2)}
\centering\scriptsize
\renewcommand{\arraystretch}{1.35}
\rowcolors{2}{iubnavy!6}{white}
\begin{adjustbox}{max width=\linewidth, max totalheight=5.6cm}
\begin{tabularx}{\textwidth}{>{\raggedright\arraybackslash}X>{\raggedright\arraybackslash}X>{\raggedright\arraybackslash}X>{\raggedright\arraybackslash}X>{\raggedright\arraybackslash}X}
  \rowcolor{iubnavy}
  \hd{Parámetro} & \hd{Símbolo} & \hd{Unidad} & \hd{Submúltiplos habituales} & \hd{Fórmula o definición}\\
  Constante de tiempo & \ensuremath{\tau} & segundo (s) & ms, µs & RC\\
  Frecuencia de corte & fc & hercio (Hz) & kHz & 1/(2\ensuremath{\pi}RC)\\
  Ángulo de fase & \ensuremath{\theta} & grado (°) & — & tan\ensuremath{^{-}}¹ de la relación reactiva/resistiva\\
  Tensión y corriente eficaces & V, I & voltio (V), amperio (A) & mV, mA & magnitud del fasor\\
\end{tabularx}
\end{adjustbox}
\vspace{2mm}
{\scriptsize Parámetros de los circuitos RC: símbolo, unidad, submúltiplos habituales y fórmula\par}
\end{frame}

% =====================================================================
\bloque{07}{Resistor en serie con un paralelo RC a 1 kHz: predecir, medir y justificar las diferencias}{50 min}{iubblue}
% =====================================================================

% ---------------------------------------------------------------------
\begin{frame}{{\normalsize Cálculo y medición de un circuito RC mixto en el laboratorio}}
\begin{columns}[c,onlytextwidth]
  \column{0.57\textwidth}
\small
\begin{itemize}\setlength{\itemsep}{4pt}
  \item La identificación y la verificación de los tres componentes cubren el criterio CE1.
  \item Siguiendo el método general: $G_2 = 1/1492 = 0.670\ \text{mS}$ y $B_C = 2\pi(1000)(97.8 \times 10^{-9}) = 0.615\ \text{mS}$, de modo que $\mathbf{Y}_p = 0.909 \angle 42.5^\circ\ \text{mS}$ y $\mathbf{Z}_p = 1100 \angle -42.5^\circ\ \Omega = 811 - j743\ \Omega$.
  \item Ambas leyes de Kirchhoff se cumplen solo en forma fasorial.
\end{itemize}
  \column{0.4\textwidth}
\begin{caja}{iubblue}{Ecuación clave}
\footnotesize \centering\small \begin{adjustbox}{max width=\linewidth}$\displaystyle \begin{gathered}\mathbf{Z}_T= 996 + (811 - j743) \\ = 1807 - j743\ \Omega \\ = 1954 \angle -22.4^\circ\ \Omega\end{gathered}$\end{adjustbox}
\end{caja}
\end{columns}
\end{frame}

% ---------------------------------------------------------------------
\begin{frame}{Circuito RC mixto R1 + (R2 \ensuremath{\parallel} C) a 1 kHz}
\centering\scriptsize
\renewcommand{\arraystretch}{1.35}
\rowcolors{2}{iubnavy!6}{white}
\begin{adjustbox}{max width=\linewidth, max totalheight=5.6cm}
\begin{tabularx}{\textwidth}{>{\raggedright\arraybackslash}X>{\raggedright\arraybackslash}X>{\raggedright\arraybackslash}X>{\raggedright\arraybackslash}X>{\raggedright\arraybackslash}X}
  \rowcolor{iubnavy}
  \hd{Magnitud} & \hd{Predicción} & \hd{Medido} & \hd{Diferencia} & \hd{Causa probable de la diferencia}\\
  I & 2.56 mA & 2.55 mA & \ensuremath{-}0.4 \% & resistencia del amperímetro\\
  VR1 & 2.55 V & 2.54 V & \ensuremath{-}0.4 \% & resolución del multímetro\\
  Vp & 2.81 V & 2.82 V & +0.4 \% & tolerancia de lectura en ca\\
  IR2 & 1.89 mA & 1.88 mA & \ensuremath{-}0.5 \% & resistencia del amperímetro\\
  IC & 1.73 mA & 1.73 mA & 0.0 \% & —\\
  \ensuremath{\Delta}t (I respecto a VS) & 62 µs & 61 µs & \ensuremath{-}1.6 \% & resolución de los cursores\\
  \ensuremath{\theta} & 22.4° & 22.0° & 0.4° & resolución de los cursores\\
\end{tabularx}
\end{adjustbox}
\vspace{2mm}
{\scriptsize Circuito RC mixto R1 + (R2 \ensuremath{\parallel} C) a 1 kHz: predicción con valores medidos y mediciones (lecturas de ejemplo)\par}
\end{frame}

% =====================================================================
\bloque{08}{Síntesis, cierre actitudinal y bibliografía}{15 min}{iuborange}
% =====================================================================

% ---------------------------------------------------------------------
\begin{frame}{Síntesis de la sesión}
\begin{columns}[T,onlytextwidth]
  \column{0.32\textwidth}
\begin{cardB}{Método general de cálculo con impedancias}
\scriptsize Pasos del análisis fasorial: reactancias, reducción, ley de Ohm y verificación con Kirchhoff
\end{cardB}
  \column{0.32\textwidth}
\begin{cardO}{Circuitos RC mixtos (serie-paralelo)}
\scriptsize Reducción de un resistor en serie con un paralelo RC y reparto de tensiones y corrientes
\end{cardO}
  \column{0.32\textwidth}
\begin{cardG}{El circuito RC como filtro}
\scriptsize Filtros pasabajas y pasaaltas, fc = 1/(2\ensuremath{\pi}RC) y el 70.7 \%
\end{cardG}
\end{columns}
\vspace{2mm}
\begin{columns}[T,onlytextwidth]
  \column{0.32\textwidth}
\begin{cardB}{Instrumentos de medición y sus efectos}
\scriptsize Resistencia interna del generador, carga de la punta del osciloscopio, ancho de banda del multímetro
\end{cardB}
  \column{0.32\textwidth}
\begin{cardO}{Parámetros, unidades, simbología}
\scriptsize Cuadro de magnitudes, unidades, prefijos, símbolos y notación fasorial
\end{cardO}
  \column{0.32\textwidth}
\begin{cardG}{Cálculo y medición de un circuito RC mixto}
\scriptsize Resistor en serie con un paralelo RC a 1 kHz: predecir, medir y justificar las diferencias
\end{cardG}
\end{columns}
\end{frame}

% ---------------------------------------------------------------------
\begin{frame}{Reflexión para llevar}
\centering
\begin{tikzpicture}
  \node[fill=iubnavy, text=white, rounded corners=6pt, text width=11.5cm, inner sep=12pt, align=center, font=\normalsize] (c) at (0,0) {\itshape «Un informe que justifica cada diferencia con una causa medible es un acto de honestidad profesional: quien lo lea, en el laboratorio o en una empresa, puede confiar en los datos sin tener que repetir las mediciones del equipo.»};
  \node[fill=iubyellow, text=iubnavy, rounded corners=3pt, font=\scriptsize\bfseries, inner sep=4pt, anchor=south west] at ([yshift=-4pt]c.north west) {Componente actitudinal};
\end{tikzpicture}
\end{frame}

% ---------------------------------------------------------------------
\begin{frame}{Referencias bibliográficas}
\small
\begin{itemize}\setlength{\itemsep}{8pt}
  \item[{$[1]$}] T. L. Floyd, Principios de circuitos eléctricos, 8.\textordfeminine{} ed. México: Pearson Educación, 2007.
  \item[{$[2]$}] R. L. Boylestad, Introducción al análisis de circuitos, 10.\textordfeminine{} ed. México: Pearson Educación, 2004.
  \item[{$[3]$}] M. F. P. Deorsola y P. Morcelle del Valle, Circuitos eléctricos. Parte 1, 1.\textordfeminine{} ed. La Plata: Editorial de la Universidad de La Plata, 2017.
  \item[{$[4]$}] Manuales del generador de funciones, del multímetro digital y del osciloscopio del laboratorio (resistencia de salida, ancho de banda y referencia de tierra); no están cargados en el wiki. (no cargada en el wiki)
\end{itemize}
\end{frame}

% ---------------------------------------------------------------------
%  CIERRE
% ---------------------------------------------------------------------
\begin{frame}[plain]
\begin{tikzpicture}[remember picture, overlay]
  \fill[iubnavy] (current page.north west) rectangle (current page.south east);
  \node[anchor=north west, text=iubyellow, font=\bfseries\fontsize{40}{44}\selectfont]
    at ([xshift=1.2cm,yshift=-1.6cm]current page.north west) {\textexclamdown{}Gracias!};
  \node[anchor=north west, text=white, font=\normalsize, text width=14cm, align=left]
    at ([xshift=1.2cm,yshift=-3.4cm]current page.north west)
    {IEL04 \textperiodcentered{} Circuitos Eléctricos I};
  \node[anchor=north west, fill=iubblue, text=white, rounded corners=1.5mm,
        font=\small\bfseries, inner sep=2.2mm, text width=11.5cm]
    at ([xshift=1.2cm,yshift=-4.4cm]current page.north west)
    {Próxima sesión \textperiodcentered{} Semana 8: Circuitos RL en serie y en paralelo: tipos e identificación de resistores e inductores, análisis con las leyes de voltajes y corrientes de Kirchhoff};
  \node[anchor=south west, text=iubyellow, font=\small, align=left]
    at ([xshift=1.2cm,yshift=0.9cm]current page.south west)
    {Ing. Sergio Martinez Castilla \quad · \quad smartinezc@unibarranquilla.edu.co};
  \fill[iubblue]   ([xshift=1.2cm,yshift=0.55cm]current page.south west) rectangle ++(1.6cm,0.15cm);
  \fill[iuborange] ([xshift=2.9cm,yshift=0.55cm]current page.south west) rectangle ++(1.6cm,0.15cm);
  \fill[iubgreen]  ([xshift=4.6cm,yshift=0.55cm]current page.south west) rectangle ++(1.6cm,0.15cm);
\end{tikzpicture}
\end{frame}

\end{document}
